\documentclass[11pt]{article}

\usepackage[margin=1in]{geometry}
\usepackage{amsmath,amssymb,amsthm,mathtools}
\usepackage{microtype}
\usepackage{listings}
\usepackage{xcolor}
\usepackage[hidelinks]{hyperref}

\newtheorem{theorem}{Theorem}
\newtheorem{lemma}[theorem]{Lemma}
\newtheorem{remark}[theorem]{Remark}
\newtheorem{definition}[theorem]{Definition}

\newcommand{\R}{\mathbb{R}}
\newcommand{\Z}{\mathbb{Z}}
\newcommand{\F}{\mathcal{F}}
\newcommand{\K}{\mathcal{K}}
\newcommand{\Jac}{\operatorname{Jac}}

\title{Reducing \textsc{3SAT} to Zero-Finding for Constant-Jacobian Polynomial Maps}
\author{Lev Stambler\thanks{This work was largely assisted by GPT-5.6 Sol,
which also generated the PDF, and was verified by Lev Stambler.}\\
\small University of Maryland, College Park\\
\small Tear Labs}
\date{}

\begin{document}

\maketitle

\begin{abstract}
We give an explicit polynomial-time reduction from \textsc{3SAT} to the
problem of deciding whether a polynomial map
\[
K:\R^N\longrightarrow\R^N
\]
has a real zero, under the promise that its Jacobian determinant is a fixed
nonzero constant.

The reduction uses a single fixed polynomial map
\[
H:\R^3\longrightarrow\R^3
\]
with three properties:
\begin{enumerate}
\item $\det JH$ is a nonzero constant;
\item the preimage $H^{-1}(0)$ consists of three explicitly known points
encoding the bits $0$ and $1$;
\item one explicitly known target $\Delta$ is omitted from the image of
$H$.
\end{enumerate}
Copies of this fixed map encode variables and clauses. The Jacobian of the
resulting map is block lower triangular, so its determinant is immediately
seen to be a nonzero constant.
\end{abstract}

\section{Statement of the reduction}

\begin{theorem}
There is a polynomial-time algorithm which, given a $3$CNF formula
$\phi$ with $n$ variables and $m$ clauses, outputs an integer-coefficient
polynomial map
\[
K_\phi:\R^{3n+3m}\longrightarrow\R^{3n+3m}
\]
such that:
\begin{enumerate}
\item
\[
\det JK_\phi\equiv (-3456)^{n+m};
\]
\item
\[
\deg K_\phi\leq 7;
\]
\item
\[
\phi\in\mathrm{SAT}
\quad\Longleftrightarrow\quad
\exists u\in\R^{3n+3m}\text{ such that }K_\phi(u)=0.
\]
\end{enumerate}
Consequently, since $3\mathrm{SAT}$ is NP-hard~\cite{Karp1972}, deciding real
zero existence for integer-coefficient polynomial maps with constant nonzero
Jacobian determinant is NP-hard.
\end{theorem}

The constant-Jacobian problem goes back to Keller's original
paper~\cite{Keller1939}.  Our proof is based on one fixed three-dimensional
gadget.  Its unscaled form is the explicit map announced by
Alpöge~\cite{Alpoge2026}; the affine rescaling below is convenient for the
reduction.

\section{The fixed polynomial gadget}

Define
\[
H=(H_1,H_2,H_3):\R^3\longrightarrow\R^3
\]
by
\begin{align}
H_1(x,y,z)
&=
12\left(
(1+xy)^3z
+
y^2(1+xy)(4+3xy)
\right)
+3,
\\
H_2(x,y,z)
&=
12\left(
y
+
3x(1+xy)^2z
+
3xy^2(4+3xy)
\right),
\\
H_3(x,y,z)
&=
12\left(
2x-3x^2y-x^3z
\right).
\end{align}

Also define the three points
\[
r_0=
\left(0,0,-\frac14\right),
\qquad
r_+=
\left(1,-\frac32,\frac{13}{2}\right),
\qquad
r_-=
\left(-1,\frac32,\frac{13}{2}\right),
\]
and the target
\[
\Delta=(7,24,8).
\]

For a target $q\in\R^3$, its preimage under $H$\footnote{In standard
terminology, this preimage is called the \emph{fiber} of $H$ over $q$.} is
\[
H^{-1}(q)=\{w\in\R^3:H(w)=q\}.
\]

\begin{lemma}[Fixed gadget]
The map $H$ satisfies:
\begin{enumerate}
\item
\[
\det JH\equiv -3456;
\]
\item
\[
H^{-1}(0)={r_0,r_+,r_-};
\]
\item
\[
H^{-1}(\Delta)=\varnothing.
\]
\end{enumerate}
\end{lemma}

\begin{proof}
The Jacobian identity is a direct polynomial calculation:
\[
\det JH=12^3(-2)=-3456.
\]

For completeness, the two preimage claims have short Gröbner-basis
certificates.  We use the standard connection between Gröbner bases,
elimination, and polynomial solution sets; see, for example,
\cite[Chs.~2--4]{CLO2025}.

With respect to lexicographic order
\[
z\succ y\succ x,
\]
a Gröbner basis for the ideal defining the preimage of zero,
\[
\langle H_1,H_2,H_3\rangle
\subseteq \mathbb{Q}[x,y,z]
\]
is
\[
\left\{
z-\frac{27}{4}x^2+\frac14,\;
y+\frac32x,\;
x^3-x
\right\}.
\]

Hence every zero satisfies
\[
x^3-x=x(x-1)(x+1)=0,
\]
so
\[
x\in\{-1,0,1\}.
\]
The other two equations then uniquely determine $y$ and $z$:
\[
y=-\frac32x,
\qquad
z=\frac{27}{4}x^2-\frac14.
\]
The resulting three points are exactly
\[
r_0,\qquad r_+,\qquad r_-.
\]

For the omitted target, a Gröbner basis for
\[
\langle
H_1-7,\;
H_2-24,\;
H_3-8
\rangle
\]
is
\[
\{1\}.
\]
Therefore this system has no solution even over $\mathbb{C}$, and in
particular has no real solution.
\end{proof}

\section{The two jobs of the fixed gadget}

The reduction uses $H$ as a small machine with two independent jobs.

First, the preimage of zero under $H$ stores one Boolean bit.  Define
\[
\beta(x,y,z)=x^2.
\]
The three possible zeroes of $H$ decode as
\[
\beta(r_0)=0,
\qquad
\beta(r_+)=\beta(r_-)=1.
\]
Consequently,
\[
H(w)=0
\quad\Longrightarrow\quad
\beta(w)\in\{0,1\}.
\]
Both bit values can be encoded: use $r_0$ for $0$, and use either $r_+$ or
$r_-$ for $1$.

Second, $H$ can hit the target $0$ but cannot hit the target $\Delta$.  This
lets us turn a Boolean value into a question about whether an equation has a
solution.

\subsection{The target selector}

For a scalar $s$, define the map $T_s:\R^3\to\R^3$ by
\[
T_s(v)=H(v)-(1-s)\Delta.
\]
The key observation is simply
\[
T_s(v)=0
\quad\Longleftrightarrow\quad
H(v)=(1-s)\Delta.
\]
Thus $s$ selects the target that $H$ is asked to hit.  For Boolean $s$ there
are only two cases:
\[
\begin{array}{c|c|c}
s & \text{target required from }H & \text{can }H\text{ hit it?} \\
\hline
1 & 0      & \text{yes} \\
0 & \Delta & \text{no}.
\end{array}
\]

\begin{lemma}[Boolean target selector]
For $s\in\{0,1\}$,
\[
T_s^{-1}(0)\neq\varnothing
\quad\Longleftrightarrow\quad
s=1.
\]
Moreover,
\[
\det D_vT_s(v)\equiv -3456.
\]
\end{lemma}

\begin{proof}
If $s=1$, then $T_1=H$, so for example $T_1(r_0)=0$.  If $s=0$, then
$T_0(v)=0$ would imply $H(v)=\Delta$, which is impossible by the fixed-gadget
lemma.  Finally, $(1-s)\Delta$ is constant with respect to $v$, so
\[
D_vT_s(v)=JH(v).
\]
The determinant claim follows from $\det JH\equiv-3456$.
\end{proof}

Nothing is claimed about non-Boolean values of $s$.  In the reduction below,
the variable equations ensure that every value fed into a target selector is
either $0$ or $1$.

\section{From a formula to a polynomial map}

Let $\phi(X_1,\dots,X_n)$ be a $3$CNF formula with clauses
$C_1,\dots,C_m$.  The construction has three steps:
\begin{enumerate}
\item force each variable block to encode a bit;
\item compute the Boolean value of every clause as a polynomial;
\item use a target selector to require every clause value to be $1$.
\end{enumerate}

\subsection{Variables}

For each variable $X_i$, introduce a block
\[
w_i=(a_i,b_i,c_i)\in\R^3
\]
and include $H(w_i)$ in the output.  At a zero of the final map this forces
$w_i\in\{r_0,r_+,r_-\}$, so
\[
X_i=\beta(w_i)
\]
defines a Boolean assignment.  Conversely, an assignment can be encoded by
choosing $w_i=r_0$ when $X_i=0$ and $w_i=r_+$ when $X_i=1$.

\subsection{Literals and clauses}

For a literal $\ell$, define its polynomial value by
\[
\operatorname{val}(\ell)=
\begin{cases}
\beta(w_i),
&
\ell=X_i,
\\[1mm]
1-\beta(w_i),
&
\ell=\neg X_i.
\end{cases}
\]

Whenever $H(w_i)=0$, each such literal value belongs to $\{0,1\}$ and agrees
with the usual Boolean evaluation under the assignment
\[
X_i=\beta(w_i).
\]

Let
\[
C_j=\ell_{j,1}\lor\ell_{j,2}\lor\ell_{j,3}
\]
be the $j$th clause. Define its satisfaction polynomial
\[
s_j(w_1,\dots,w_n)=1-
\prod_{k=1}^{3}
\left(
1-\operatorname{val}(\ell_{j,k})
\right).
\]

On Boolean inputs,
\[
s_j\in\{0,1\},
\]
and
\[
s_j=1
\quad\Longleftrightarrow\quad
C_j\text{ is satisfied}.
\]

For each clause, introduce a fresh witness block
\[
v_j\in\R^3
\]
and include the output
\[
T_{s_j}(v_j)=H(v_j)-(1-s_j)\Delta.
\]

Because $s_j\in\{0,1\}$ at every solution of the variable equations, the
target-selector lemma implies
\[
\exists v_j:\ T_{s_j}(v_j)=0
\quad\Longleftrightarrow\quad
s_j=1.
\]

In other words, $v_j$ is a witness that exists exactly when clause $C_j$ is
satisfied.

\subsection{The final map}

Let $\phi$ have $n$ variables and $m$ clauses. Define
\[
K_\phi:
\R^{3n+3m}
\longrightarrow
\R^{3n+3m}
\]
by
\[
K_\phi
\left(
w_1,\dots,w_n,
v_1,\dots,v_m
\right)=
\begin{pmatrix}
H(w_1)
\\
\vdots
\\
H(w_n)
\\
T_{s_1}(v_1)
\\
\vdots
\\
T_{s_m}(v_m)
\end{pmatrix}.
\]

Equivalently,
\[
K_\phi=
\begin{pmatrix}
H(w_1)
\\
\vdots
\\
H(w_n)
\\
H(v_1)-(1-s_1(w))\Delta
\\
\vdots
\\
H(v_m)-(1-s_m(w))\Delta
\end{pmatrix}.
\]

\section{Why the reduction works}

The construction was designed so that its zeroes and satisfying assignments
are the same information in two different forms.

\begin{lemma}[Correctness]
The map $K_\phi$ has a real zero if and only if $\phi$ is satisfiable.
\end{lemma}

\begin{proof}
Suppose first that $K_\phi(w,v)=0$.  Every variable block then satisfies
$H(w_i)=0$, so the values $X_i=\beta(w_i)$ form a Boolean assignment.  Every
clause block satisfies
\[
T_{s_j}(v_j)=0.
\]
Since $s_j$ is Boolean, the target-selector lemma forces $s_j=1$.  Hence every
clause is true under the decoded assignment, and $\phi$ is satisfiable.

Conversely, suppose that $\phi$ has a satisfying assignment $\sigma$.  Encode
it by choosing
\[
w_i=
\begin{cases}
r_0,&\sigma(X_i)=0,\\
r_+,&\sigma(X_i)=1.
\end{cases}
\]
Then $H(w_i)=0$ and $\beta(w_i)=\sigma(X_i)$.  Since $\sigma$ satisfies every
clause, every $s_j=1$.  Choose $v_j=r_0$ for all $j$; then
\[
T_{s_j}(v_j)=T_1(r_0)=H(r_0)=0.
\]
Thus every output coordinate of $K_\phi$ vanishes.  We have proved
\[
K_\phi^{-1}(0)\neq\varnothing
\quad\Longleftrightarrow\quad
\phi\in\mathrm{SAT}.
\]
\end{proof}

\section{Completing the proof of the reduction}

It remains to check that the construction has the promised Jacobian,
degree, coefficient, and complexity bounds.

\begin{lemma}[Constant Jacobian]
The Jacobian determinant of $K_\phi$ is the nonzero constant
\[
\det JK_\phi\equiv(-3456)^{n+m}.
\]
\end{lemma}

\begin{proof}

Order the input blocks as
\[
w_1,\dots,w_n,v_1,\dots,v_m.
\]

The intuition is that the clause gadgets may look backward at the encoded
variables, but no variable gadget looks forward at the clause witnesses.
Algebraically, the variable outputs depend only on the $w$-blocks.  The
$j$th clause output may depend on all $w$-blocks through $s_j$, but among the
witness blocks it depends only on $v_j$.

Therefore the Jacobian has block lower-triangular form
\[
JK_\phi=
\begin{pmatrix}
D_{\mathrm{var}} & 0
\\
B & D_{\mathrm{clause}}
\end{pmatrix},
\]
where
\[
D_{\mathrm{var}}=\operatorname{diag}
\left(
JH(w_1),\dots,JH(w_n)
\right)
\]
and
\[
D_{\mathrm{clause}}=\operatorname{diag}
\left(
JH(v_1),\dots,JH(v_m)
\right).
\]

The lower-left block $B$ contains derivatives of the $s_j$ with respect to
the variable blocks.  Its exact form is irrelevant to the determinant.  The
diagonal clause blocks are still $JH(v_j)$ because the target shift
$(1-s_j)\Delta$ does not depend on $v_j$.

Since the determinant of a block triangular matrix is the product of the
determinants of its diagonal blocks,
\begin{align}
\det JK_\phi
&=
\det D_{\mathrm{var}}
\det D_{\mathrm{clause}}
\\
&=
\prod_{i=1}^{n}\det JH(w_i)
\prod_{j=1}^{m}\det JH(v_j)
\\
&=
(-3456)^n(-3456)^m
\\
&=
(-3456)^{n+m}.
\end{align}

Thus the Jacobian determinant is a nonzero constant polynomial.
\end{proof}

\begin{lemma}[Degree, coefficients, and size]
The map $K_\phi$ has integer coefficients, degree at most $7$, and a
description that can be constructed in time polynomial in the size of
$\phi$.
\end{lemma}

\begin{proof}

The map $H$ has degree $7$, while $\beta$ has degree $2$.  Each $s_j$ is a
product of at most three degree-$2$ literal expressions, so
\[
\deg s_j\leq 6.
\]

Each clause output has the form
\[
H(v_j)-(1-s_j)\Delta.
\]
Therefore
\[
\deg K_\phi
\leq
\max\{7,6\}=7.
\]

All coefficients are integers:
\begin{itemize}
\item $H$ has integer coefficients;
\item $\Delta=(7,24,8)$ is integral;
\item every $s_j$ has integer coefficients.
\end{itemize}

There are exactly
\[
3n+3m
\]
input and output coordinates. Since each variable and clause contributes only
a constant-size polynomial block, the description of $K_\phi$ has size
polynomial in the description length of $\phi$.
\end{proof}

\begin{proof}[Proof of Theorem~1]
The correctness lemma gives
\[
K_\phi^{-1}(0)\neq\varnothing
\quad\Longleftrightarrow\quad
\phi\in\mathrm{SAT}.
\]
The two preceding lemmas give the claimed Jacobian, degree, coefficient, and
polynomial-time construction bounds.  Since $3\mathrm{SAT}$ is
NP-hard~\cite{Karp1972}, the zero-existence problem in the theorem is NP-hard.
\end{proof}

\appendix

\section{Exact symbolic verification}

The only special algebraic facts are:
\[
\det JH=-3456,
\qquad
H^{-1}(0)={r_0,r_+,r_-},
\qquad
H^{-1}(\Delta)=\varnothing.
\]

They can be verified exactly over $\mathbb{Q}$ with the following script.

\begin{lstlisting}[language=Python,basicstyle=\ttfamily\small]
import sympy as sp

x, y, z = sp.symbols("x y z")
u = 1 + x*y

P = u**3*z + y**2*u*(4 + 3*x*y)
Q = y + 3*x*u**2*z + 3*x*y**2*(4 + 3*x*y)
R = 2*x - 3*x**2*y - x**3*z

F = sp.Matrix([P, Q, R])

# The unscaled map has constant Jacobian -2.

assert sp.factor(F.jacobian((x, y, z)).det()) == -2

# H = 12(F - (-1/4,0,0)).

H = sp.Matrix([
12*P + 3,
12*Q,
12*R,
])

assert sp.factor(H.jacobian((x, y, z)).det()) == -3456

# Preimage of zero under H:

# equivalent to F = (-1/4,0,0).

G_zero = sp.groebner(
[
P + sp.Rational(1, 4),
Q,
R,
],
z, y, x,
order="lex",
)

print(G_zero)

# z - 27*x**2/4 + 1/4

# y + 3*x/2

# x**3 - x

# Preimage of (7,24,8) under H:

# equivalent to F = (1/3,2,2/3).

G_omitted = sp.groebner(
[
P - sp.Rational(1, 3),
Q - 2,
R - sp.Rational(2, 3),
],
z, y, x,
order="lex",
)

print(G_omitted)

# 1

\end{lstlisting}

\begin{thebibliography}{9}

\bibitem{Alpoge2026}
L.~Alpöge,
``Announcement of an explicit three-variable counterexample to the Jacobian
conjecture,'' X post, 19 July 2026.
\url{https://x.com/__alpoge__/status/2079028340955197566}.

\bibitem{CLO2025}
D.~A. Cox, J.~Little, and D.~O'Shea,
\emph{Ideals, Varieties, and Algorithms: An Introduction to Computational
Algebraic Geometry and Commutative Algebra}, 5th ed., Springer, 2025.
\url{https://doi.org/10.1007/978-3-031-91841-4}.

\bibitem{Karp1972}
R.~M. Karp,
``Reducibility among combinatorial problems,'' in
\emph{Complexity of Computer Computations}, Plenum Press, 1972, pp.~85--103.
\url{https://doi.org/10.1007/978-1-4684-2001-2_9}.

\bibitem{Keller1939}
O.-H. Keller,
``Ganze Cremona-Transformationen,''
\emph{Monatshefte für Mathematik und Physik}, vol.~47, pp.~299--306, 1939.
\url{https://doi.org/10.1007/BF01695502}.

\end{thebibliography}

\end{document}
